\section{Case Study: $n$-gram Performance on InstructCoder}
\label{subsect:ngram-bleu-score}


The $n$-gram SD is especially attractive for speculative decoding because it is training-free.
We observe that $n$-gram achieves particularly strong speedups on the code-editing workload, such as InstructCoder, and we therefore conduct a focused case study to understand the underlying reasons.

We hypothesize that its advantage stems from the \emph{local repetition} inherent in code-editing tasks: when upcoming tokens have already appeared in the prompt, $n$-gram lookup is more likely to propose correct continuations that are subsequently accepted by the target model during verification.
To test this hypothesis, we quantify prompt–output overlap using BLEU-$n$~\cite{papineni2002bleu}.
BLEU-$n$ measures the fraction of overlapping $n$-grams, with higher values indicating stronger local reuse or copying.
For each request, we compute the BLEU score between the prompt and the output generated without speculative decoding, and then group them into 5 intervals: requests with BLEU score of 0-0.2, 0.2-0.4, 0.4-0.6, 0.6-0.8 and 0.8-1.0.
Within each interval, we compare the speedups achieved by $n$-gram and by EAGLE/EAGLE-3.
We report the results using BLEU-4 score by default~\cite{papineni2002bleu}. In our experiments, BLEU-4 gives the clearest speedup heatmap and shows a clear cutoff where, above a certain overlap range, $n$-gram beats EAGLE/EAGLE-3 for every batch size. We have also repeated the analysis with BLEU-1 through BLEU-10 scores, and observe the same overall patterns.

\begin{quoting}[leftmargin=1.5em, rightmargin=1.5em, vskip=0pt]
\textbf{Overlap Drives $n$-gram Speedups.}
\emph{$n$-gram works best on code editing because the output often reuses short spans that already appear in the prompt. As overlap increases, $n$-gram speedup rises and eventually surpasses EAGLE/EAGLE-3.}
\end{quoting}
We observe that higher BLEU-$n$ scores strongly correlate with greater $n$-gram speedup (\autoref{fig:ngram_bleu_heatmap}). 
For requests with lower BLEU-$n$ scores (little overlap), $n$-gram underperforms due to inaccurate proposals. Once the BLEU-$n$ score exceeds a certain threshold, (e.g. larger than $0.6$ for Llama3.1-8B), $n$-gram consistently outperforms EAGLE and EAGLE-3 across all batch sizes. With a proposal length of 3, it achieves up to $53\%$ higher speedup than EAGLE and EAGLE-3. When we further increase the proposal length of $n$-gram to 5, the boundary where $n$-gram performs better is clearer, and this brings an even higher performance of up to $100\%$ higher speedup than EAGLE/EAGLE-3.
This is likely because a larger proposal length allows $n$-gram to reuse longer repeated spans in the prompt or recent context when the overlap is high. 
Overall, these results confirm that $n$-gram speculation benefits directly from prompt-level repetition in code-editing workloads. 

Finally, our BLEU analysis measures overlap only between the complete prompt and the complete generated output. In practice, $n$-gram can match and reuse tokens from the full in-context history, which includes both the prompt and the tokens generated so far. We leave an analysis that measures overlap with respect to this evolving context for future work.

% A consistent trend emerges for non-reasoning models (i.e. Llama3.1-8B, Llama3-70B) and non-reasoning datasets (i.e. InstructCoder, ShareGPT and CNNDailyMail): as BLEU increases, the speedup of $n$-gram speculation rises and eventually surpasses that of the Eagle/Eagle-3 methods. Figure~X (Llama3.1-8B and Llama3-70B on \texttt{InstructCoder}) illustrates this relationship, confirming that greater prompt--output overlap directly results in better $n$-gram efficiency. \jerry{Add Bar chart here!} 

% However, this does not hold true for reasoning models (i.e. Qwen3-8B) and reasoning datasets (i.e. AIME22-24, GPQA-MAIN). We found that the similarity between the prompt and generated outputs are often too low for 

% \paragraph{Cutoffs where $n$-gram beats Eagle/Eagle-3.}
% To quantify when $n$-gram speculation becomes more effective than Eagle/Eagle-3 methods, we define a \emph{BLEU score cutoff} for each (model, dataset) pair. The cutoff is the smallest BLEU score bucket where the speedup of $n$-gram exceeds that of the compared Eagle or Eagle-3 method and remains higher in the subsequent bucket. We sweep and compute the BLEU-1 to BLEU-10 scores for every request across all configurations and identify the most appropriate cutoffs among them.
% \autoref{tab:ngram-cutoffs} summarizes the observed cutoffs across all configurations.

% \begin{table}[t]
% \centering
% \caption{BLEU score cutoffs where $n$-gram exceeds other speculative decoding (SD) methods.
% “Cmp.” denotes the SD method being compared. “Eg/Eg-3” refers to Eagle/Eagle-3.
% “Weak” indicates that fewer than 10 requests fall into BLEU buckets higher than the cutoff.
% "NF" means no reliable cutoff found because average speedups do not cross. 
% "NF-D" means majority of the requests have low BLEU score, "NF-C" means no reliable cutoff found despite that requests have a range of BLEU scores.
% \jerry{Double check the results for Qwen3-8B.}}
% \label{tab:ngram-cutoffs}
% \resizebox{\linewidth}{!}{
% \begin{tabular}{l l c c}
% \toprule
% \textbf{Model} & \textbf{Dataset} & \textbf{Cmp.} & \textbf{Cutoff} \\
% \midrule
% Llama3.1-8B & InstructCoder & Eg/Eg-3   & BLEU-4 $> 0.5$ \\
%             & ShareGPT      & Eg/Eg-3   & BLEU-4 $> 0.6$ (weak) \\
%             & CNNDailyMail  & Eg/Eg-3   & BLEU-4 $> 0.4$ (weak) \\
%             & GSM8K         & Eg/Eg-3 & NF-D(?) \\
% \midrule
% Llama3-70B  & InstructCoder & Eg   & BLEU-4 $> 0.6$ \\
%             & ShareGPT      & Eg   & BLEU-3 $> 0.7$ (weak) \\
%             & CNNDailyMail  & Eg   & NF-C \\
%             & GSM8K         & Eg   & NF-C \\
% \midrule
% Qwen3-8B    & InstructCoder & Eg-3   &  BLEU-3 $> 0.35$ (weak)\\
%             & ShareGPT      & Eg-3   &  NF-C\\
%             & CNNDailyMail  & Eg-3   &  NF-C\\
%             & GSM8K         & Eg-3   &  NF-D\\
%             & AIME21-24     & Eg-3   &  NF-D\\
%             & GPQA-Main     & Eg-3   &  NF-D\\
% \bottomrule
% \end{tabular}
% }
% \end{table}

